\begin{lemma}
\label{lemma-Noetherian-power}
\begin{slogan}
An ideal in a Noetherian ring is nilpotent if each element
of the ideal is nilpotent.
\end{slogan}
Let $R$ be a Noetherian ring. Let $I, J$ be ideals of $R$.
Suppose $J \subset \sqrt{I}$. Then $J^n \subset I$ for some $n$.
In particular, in a Noetherian ring the notions of
``locally nilpotent ideal''
and ``nilpotent ideal'' coincide.
More generally, the conclusion $J^n\subset I$ needs only finite
generation of $J$, not Noetherianity of $R$. If
$J=(f_1,\ldots,f_s)$ and $f_i^{d_i}\in I$ with $d_i\geq1$, one may
take $n=1+\sum_{i=1}^s(d_i-1)$.
\end{lemma}

\begin{proof}
Choose generators $J=(f_1,\ldots,f_s)$ and positive exponents
$d_i$ such that $f_i^{d_i}\in I$. Put $N=1+\sum_{i=1}^s(d_i-1)$.
The ideal $J^N$ is generated by the full list of monomials
$\prod_i f_i^{\alpha_i}$ with $\alpha_i\geq0$ and $\sum_i\alpha_i=N$.
Every such list has $\alpha_i\geq d_i$ for some $i$, so that monomial
belongs to $I$. Hence $J^N\subset I$. For $s=0$, $J=0$ and $N=1$
works. The original larger choice $d_1+\cdots+d_s+1$ also works,
since every higher power of $J$ is contained in $J^N$.
Only the existence of a finite generating set for $J$ was used.

\medskip\noindent
The bound cannot be decreased using only the displayed exponents.
In $B=\mathbf Z[T_1,\ldots,T_s]/(T_1^{d_1},\ldots,T_s^{d_s})$,
the residue classes of monomials with $0\leq\alpha_i<d_i$ form a
$\mathbf Z$-basis: the defining monomial ideal is the free subgroup
spanned by all the other monomials. Therefore
$\prod_i T_i^{d_i-1}$ is a nonzero element of
$(T_1,\ldots,T_s)^{N-1}$. Together with the upper bound this proves
sharpness, including the empty product in degree zero.
\end{proof}